\documentclass[10pt,a4paper]{amsart}
\usepackage[margin=19mm]{geometry}
\usepackage{amsmath,amssymb,mathtools}
\usepackage[scr=rsfs,cal=euler]{mathalfa}
\usepackage{xfrac}
\usepackage[unicode,hidelinks,bookmarks=false]{hyperref}
\newcommand{\ie}{i.e.\ }
\newcommand{\cf}{cf.\ }
\newcommand{\resp}{respectively\ }
\newcommand{\Subsection}{\subsection}
\providecommand{\nobreakdash}{\nobreak\hbox{-}}
\setlength{\emergencystretch}{2em}
\multlinegap0pt
\usepackage{bm}
\usepackage{bbm}
\usepackage{dsfont}
\usepackage{xspace}
\usepackage{cancel}
\usepackage{mathtools}
\usepackage{amsthm}
\makeatletter
\@ifundefined{theorem}{\newtheorem{theorem}{Theorem}}{}
\@ifundefined{assumption}{\newtheorem{assumption}[theorem]{Assumption}}{}
\@ifundefined{lemma}{\newtheorem{lemma}[theorem]{Lemma}}{}
\@ifundefined{proposition}{\newtheorem{proposition}[theorem]{Proposition}}{}
\@ifundefined{remark}{\newtheorem{remark}[theorem]{Remark}}{}
\makeatother
\newcommand{\bx}{\mathbf{x}}
\renewcommand{\Im}{\operatorname{Im}}
\renewcommand{\Re}{\operatorname{Re}}
\renewcommand{\vec}[1]{\mathbf{#1}}
\title[Limiting absorption principle for a hybrid resonance in a tw]{Limiting absorption principle for a hybrid resonance in a two-dimensional cold plasma}
\author{Maryna Kachanovska, \'Etienne Peillon}
\date{Source: arXiv:2510.15153v1 (CC BY 4.0)}
\begin{document}
\maketitle

\noindent\emph{Source and licence.} Limiting absorption principle for a hybrid resonance in a two-dimensional cold plasma. Version arXiv:2510.15153v1 (CC BY 4.0), distributed under the Creative Commons Attribution 4.0 International licence, as recorded in the arXiv licence metadata for this exact version and shown on the abstract page https://arxiv.org/abs/2510.15153v1. Full rights notice: licenses/arxiv-2510.15153-CC-BY-4.0.txt.\par\smallskip

\noindent\textit{Editorial scope.} Counted unit: the Proposition whose source label is \texttt{proposition:jbound\_basic} in the article (source0.tex lines 3502--3516, source label \texttt{proposition:jbound\_basic}), delivered together with the article's own complete original proof of it (source0.tex lines 3517--3659, one \texttt{proof} environment of 143 lines). No mathematics is written, repaired or re-derived: statement and proof are the article's own text line for line. Required context retained uncounted: assump:matrices (lines 293--326); eq:unu\_orig (lines 1156--1163); eq:est\_main\_v2 (lines 1306--1310); eq:childelta (lines 1369--1377); eq:important\_mappings (lines 1391--1398); eq:Ju (lines 1409--1412); eq:Unudelta (lines 1425--1430); eq:stabU (lines 1437--1440); eq:stabFdelta (lines 1442--1445); lem:ellreg (lines 1447--1450); lem:plancherel (lines 3403--3407); eq:zero\_trace (lines 3414--3417); lem:commutator (lines 3419--3429); rem:regularity (lines 3430--3433); lem:self\_adjoint\_commutator (lines 3483--3487). Every definition, display and label of those lines is retained unchanged. Omitted from this package and not redistributed: 1--292, 327--1155, 1164--1305, 1311--1368, 1378--1390, 1399--1408, 1413--1424, 1431--1436, 1441--1441, 1446--1446, 1451--3402, 3408--3413, 3418--3418, 3434--3482, 3488--3501, 3660--4390. Nothing inside a retained excerpt is omitted or altered: every retained body is delivered exactly as the article's lines are written, so no omission receipt of any kind accompanies this package. Citations: the retained excerpts carry no citation command at all, so no bibliography entry of the article is transcribed and no reference is redistributed. Reference adaptation: none was needed and none was made; every reference inside the delivered unit resolves inside it, so no unresolved reference or citation is produced. Printed slips kept literal: no printed word, symbol, hypothesis or number of the counted statement or of its proof was altered. Layout and class: the article's own class is \texttt{elsarticle} with packages including amsmath, amsthm, amsfonts, ulem, amssymb, xcolor, bbold, enumitem, tikz, tabularx, appendix, geometry, mathabx, mathbbol; the wrapper is the lane's pinned \texttt{amsart} editorial wrapper at 10pt on A4, which restates the theorem-like environments the retained text needs and adds the article's own macro definitions that the retained text needs (\texttt{\textbackslash Im}, \texttt{\textbackslash Re}, \texttt{\textbackslash bx}, \texttt{\textbackslash vec}), with extra wrapper packages (bm, bbm, dsfont, xspace, cancel, mathtools). The delivered numbering is the unit's own. Assets: no retained span contains a figure environment, a graphics-inclusion command, a PDF-page inclusion, a table or a TikZ picture; no image file is copied into this package, the package declares no asset dependency and no image file is redistributed with it. Licence: version arXiv:2510.15153v1 of this article is distributed under the Creative Commons Attribution 4.0 International licence, as recorded in the arXiv licence metadata for this exact version and shown on the abstract page https://arxiv.org/abs/2510.15153v1. A full rights notice is delivered at licenses/arxiv-2510.15153-CC-BY-4.0.txt. Characters: every retained excerpt is pure ASCII, so the delivered engine, XeLaTeX with its default Latin Modern text font, reports no missing character.
\begin{assumption}
	\label{assump:matrices}
	\begin{itemize}
		\item $\mathbb{A}, \, \mathbb{T}\in {C}^{1,1}(\overline{\Omega}; \mathbb{C}^{2\times 2})$. 
		\item For all $\bx\in \overline{\Omega}$, $\mathbb{A}(\bx), \mathbb{T}(\bx)$ are both Hermitian, positive-definite matrices. 
		
		
		In particular,
		defining for $\vec{p}, \vec{v}\in \mathbb{C}^2$, 
		\begin{align*}
			\vec{p}\cdot \vec{v}=p_1v_1+p_2v_2,\quad |\vec{p}|^2=\|\vec{p}\|_{\mathbb{C}^2}^2=\vec{p}\cdot\overline{\vec{p}},
		\end{align*}
		it holds, for all $\vec{p}\in \mathbb{C}^2$, all $\vec{x}\in \overline{\Omega}$,  $
		\mathbb{T}(\bx)\vec{p}\cdot\overline{\vec{p}}\geq c_{\mathbb{T}}\|\vec{p}\|^2_{\mathbb{C}^2}, \quad 			\mathbb{A}(\bx)\vec{p}\cdot\overline{\vec{p}}\geq c_{\mathbb{A}}\|\vec{p}\|^2_{\mathbb{C}^2}, $ where $c_{\mathbb{T}}, c_{\mathbb{A}}>0$. 
		\item  Moreover, $\mathbb{A}$ and $\mathbb{T}$ satisfy the following periodicity constraints: 
		\begin{align*}
			\partial_y^k\mathbb{A}(.,\ell)=	\partial_y^k\mathbb{A}(.,-\ell), \qquad 	\partial_y^k\mathbb{T}(.,\ell)=	\partial_y^k\mathbb{T}(.,-\ell), \quad k=0,\,1. 
		\end{align*}
	\end{itemize}
	We will use the following notation for the values of $\mathbb{A}$ and $\mathbb{T}$ on $\Sigma$:
	\begin{align*}
		\left.	\mathbb{A}\right|_{\Sigma}=\left(
		\begin{matrix}
			a_{11} & a_{12}\\
			\overline{a_{12}} & a_{22}
		\end{matrix}
		\right), \qquad \left.	\mathbb{T}\right|_{\Sigma}=\left(
		\begin{matrix}
			t_{11} & t_{12}\\
			\overline{t_{12}} & t_{22}
		\end{matrix}
		\right).
	\end{align*}
\end{assumption}

\begin{align}
	\label{eq:unu_orig}
	\begin{split}
		& \operatorname{div}(( x \mathbb{A}+i\nu\mathbb{T})\nabla u^{\nu})=f \text{ in }\Omega, \\
		&\gamma_0^{\Gamma_p\cup\Gamma_n}u^{\nu}=0, \text{ periodic BCs at }\Gamma^{\pm}.
	\end{split}
	\tag{P$\nu$}
\end{align}

\begin{align}
	\label{eq:est_main_v2}
	&\|u^{\nu}\|_{L^2(\Omega)}+\|x\nabla u^{\nu}\|_{L^2(\Omega)}+\nu^{1/2}\|\nabla u^{\nu}\|_{L^2(\Omega)}\lesssim \|f\|_{L^2(\Omega)}.
	%\tag{StH1}
\end{align}

\begin{align}
	\label{eq:childelta}
	\mathbb{R}^2\ni \vec{x}=(x,y)\mapsto	\chi_{\ell,\delta}(x,y):=\left\{
	\begin{array}{ll}
		1, & |y|\leq \ell+\delta/2, \\
		0,& |y|>\ell+\delta, \\
		\in (0,1), & \text{ otherwise. }
	\end{array}	\right.
\end{align}

\begin{align}
	\label{eq:important_mappings}
	\begin{split}
		&\|u\|_{\Omega}\leq \|\mathcal{E}_{\delta} u\|_{\mathbb{R}^2_a}\leq C_{\delta}\|u\|_{\Omega}, \quad \|\partial_y u\|_{\Omega}\leq \|\partial_y \mathcal{E}_{\delta} u\|_{\mathbb{R}^2_a}\leq C_{\delta}(\|\partial_y u\|_{\Omega}+\|u\|_{\Omega}),\\
		&\|\partial_y (\mathbb{P}\nabla u)\|_{\Omega}\leq \|\partial_y(\mathbb{P} \nabla\mathcal{E}_{\delta}u)\|_{\mathbb{R}^2_a}\leq C_{\delta,\mathbb{P}}( \|\mathbb{P}\partial_y\nabla u\|_{\Omega}+\|\partial_y u\|_{\Omega}+\|u\|_{\Omega}),\quad \mathbb{P}\in \mathbb{R}^{2\times 2},\\
		&\|x\nabla \mathcal{E}_{\delta}u\|_{\mathbb{R}^2_a}\leq C_{\delta} \|u\|_{\mathcal{V}_{sing}(\Omega)}.
	\end{split}
\end{align} 

\begin{align}
	\label{eq:Ju}
	\|\mathcal{J}^2 u\|^2_{L^2(\mathbb{R}^2_a)}=\int_{-a}^a\int_{\mathbb{R}}(1+\xi_y^2)|\mathcal{F}_yu(x,\xi_y)|^2d\xi_y\, dx=\|u\|^2_{L^2(\mathbb{R}^2_a)}+\|\partial_y u\|^2_{L^2(\mathbb{R}^2_a)},\quad \forall u\in H^1(\mathbb{R}^2_a).
\end{align}

\begin{align}
	\label{eq:Unudelta}
	\operatorname{div}\left((x\mathbb{A}+i\nu\mathbb{T})\nabla U^{\nu}_{\delta}\right)&=F^{\nu}_{\delta} \text{ in }\mathbb{R}^2_a, \\
	\nonumber
	F^{\nu}_{\delta}&=\mathcal{E}_{\delta}f+U^{\nu}\operatorname{div}((x\mathbb{A}+i\nu\mathbb{T})\nabla \chi_{\ell,\delta})+\nabla U^{\nu}\cdot (x\mathbb{A}+i\nu\mathbb{T})\nabla \chi_{\ell,\delta}+\nabla \chi_{\ell,\delta}\cdot (x\mathbb{A}+i\nu\mathbb{T})\nabla U^{\nu}.
\end{align}

\begin{align}
	\label{eq:stabU}
	\|U^{\nu}_{2\delta}\|_{L^2(\mathbb{R}^2_a)}+\|x\nabla U^{\nu}_{2\delta}\|_{L^2(\mathbb{R}^2_a)}+\nu^{1/2}\|\nabla U^{\nu}_{2\delta}\|_{L^2(\mathbb{R}^2_a)}&\leq   c_{2,\delta}\|f\|_{L^2(\Omega)},
\end{align}

\begin{align}
	\label{eq:stabFdelta}
	\|F^{\nu}_{\delta}\|_{L^2(\mathbb{R}^2_a)}\leq c_{3,\delta} \|f\|_{L^2(\Omega)}, \quad \forall 0<\nu<1. 
\end{align}

\begin{lemma}
	\label{lem:ellreg}
	For all $\nu>0$, $U^{\nu}_{\delta}\in H^2(\mathbb{R}^2_a)$; also $\gamma_{0}^{\partial\mathbb{R}^2_a}\left(\partial_y U^{\nu}_{\delta}\right)=0$, $\lambda\in \{n,p\}$.
\end{lemma}

\begin{lemma}
	\label{lem:plancherel}
	For all $u, v\in H^1(\mathbb{R}^2_a)$, it holds that $
		(\mathcal{J}u,v)_{L^2(\mathbb{R}^2_a)}=(u,\mathcal{J}v)_{L^2(\mathbb{R}^2_a)}.	$
\end{lemma}

\begin{align}
	\label{eq:zero_trace}
	\gamma_0^{x=\pm a} \mathcal{J} u=0,\qquad \gamma_0^{x=\pm a} \mathcal{J}^2 u=0.
\end{align}

\begin{lemma}
	\label{lem:commutator}
	Let $n\in \{1,2\}$, $\beta_0\in H^{2}(\mathbb{R}^2_a)$. Then, there exists $C>0$, s.t. for all $p\in L^2(\mathbb{R}_a^2)$, $n\in \{1,2\}$,
	\begin{align*}
		&\|[\mathcal{J}^n,\beta_0]p\|_{L^2(\mathbb{R}^2_a)}\leq C \|p\|_{L^2(\mathbb{R}^2_a)}.
	\end{align*}
%	Similarly, for all $p\in L^2_1(\mathbb{R}_a^2)$, $n\in \{1,2\}$,  
%	\begin{align*}
%		\|[\mathcal{J}^{n},x\beta_0]p\|_{L^2(\mathbb{R}^2_a)}\leq C \|p\|_{L^2_1(\mathbb{R}^2_a)}.
%	\end{align*}
\end{lemma}

\begin{remark}
	\label{rem:regularity}
We will often apply the above result in the case when $\beta_0\in C^{1,1}([-a,a]\times\mathbb{R}; \mathbb{C})$ and is compactly supported.
\end{remark}

\begin{lemma}
	\label{lem:self_adjoint_commutator}
	Let $E:=(L^2(\mathbb{R}^2_a))^2$. 
	Let $\mathcal{A}: D(\mathcal{A})\rightarrow E$ be a self-adjoint operator, with a domain $D(\mathcal{A})\subset E$. Then, for all $\vec{v}\in E$ s.t. $\mathcal{J}\vec{v}, \mathcal{J}^2\vec{v}\in D(\mathcal{A}) \text{  and   }\mathcal{A}\vec{v}\in \left(H^1(\mathbb{R}^2_a)\right)^2,$ it holds that $$\Im ([\mathcal{J}\mathbb{I}, \mathcal{A}]\vec{v},\mathcal{J} \vec{v})_{E}=\frac{1}{2i}([\mathcal{J}^2\mathbb{I}, \mathcal{A}]\vec{v},\vec{v})_{E}.$$ 
\end{lemma}

\begin{proposition}
	\label{proposition:jbound_basic}
	Let $(u^{\nu})_{\nu>0}$ solve \eqref{eq:unu_orig},  and $U^{\nu}_{\delta}=\mathcal{E}_{\delta}u^{\nu}$ satisfy \eqref{eq:Unudelta}. Then there exists $C>0$, s.t.  
	\begin{align}
		\label{eq:lemboundnu}
		\nu\|\mathcal{J}\nabla U^{\nu}_{\delta}\|^2\leq C \left(\|f\|^2+
		\|f\|\|\partial_y u^{\nu}\|\right), \text{ for all }0<\nu<1.
	\end{align}
	%	Also, 
	%	\begin{align}
		%		\label{eq:lemboundnu2}
		%		\nu\|\mathcal{J}\nabla U^{\nu}_{\delta})\|^2\lesssim  \|f\|^2+
		%		\|f\|\|\partial_y u^{\nu}\|.
		%	\end{align}
\end{proposition}

\begin{proof}
To prove \eqref{eq:lemboundnu}, we rewrite the extended problem \eqref{eq:Unudelta} in a more convenient form. In particular, we decompose the matrix $\mathbb{A}=\mathbb{D}+\mathbb{H}$, with $\mathbb{D}=\operatorname{diag}\mathbb{A}=\operatorname{diag}(\mathbb{A}_{11}, \mathbb{A}_{22})$. Next, we define the modified matrices:
\begin{align}
	\label{eq:ptilde}
	\begin{split}
		\widetilde{\mathbb{A}}&:=\mathbb{A}_{11}^{-1}\mathbb{A}\chi_{\ell,2\delta}+\mathbb{I}(1-\chi_{\ell,2\delta})=\mathbb{I}+\mathbb{E}_{\mathbb{A}}, \quad \mathbb{E}_{\mathbb{A}}=(\mathbb{A}_{11}^{-1}\mathbb{A}-\mathbb{I})\chi_{\ell,2\delta},\\
		\widetilde{\mathbb{T}}&:=\mathbb{A}_{11}^{-1}\mathbb{T}\chi_{\ell,2\delta}+\mathbb{I}(1-\chi_{\ell,2\delta})=\mathbb{I}+\mathbb{E}_{\mathbb{T}}, \quad \mathbb{E}_{\mathbb{T}}=(\mathbb{A}_{11}^{-1}\mathbb{A}-\mathbb{I})\chi_{\ell,2\delta},\\
		\widetilde{\mathbb{D}}&:=\mathbb{A}_{11}^{-1}\mathbb{D}\chi_{\ell,2\delta}+\mathbb{I}(1-\chi_{\ell,2\delta})=\mathbb{I}+\mathbb{E}_{\mathbb{D}}, \quad \mathbb{E}_{\mathbb{D}}=(\mathbb{A}_{11}^{-1}\mathbb{D}-\mathbb{I})\chi_{\ell,2\delta}, \\
		\widetilde{\mathbb{H}}&:=\mathbb{A}_{11}^{-1}\mathbb{H}\chi_{\ell,2\delta}+\mathbb{I}(1-\chi_{\ell,2\delta})=\mathbb{I}+\mathbb{E}_{\mathbb{H}}, \quad \mathbb{E}_{\mathbb{H}}=(\mathbb{A}_{11}^{-1}\mathbb{H}-\mathbb{I})\chi_{\ell,2\delta}
	\end{split}
\end{align}  
In the above, $\chi_{\ell,2\delta}$ is the same truncation function as in \eqref{eq:childelta}. The above defined matrix-valued functions are constant and equal to $\mathbb{I}$ for $|y|\geq \ell+2\delta$. The matrices $\widetilde{\mathbb{A}}, \widetilde{\mathbb{T}}, \widetilde{\mathbb{D}}, \widetilde{\mathbb{H}}$ are Hermitian and positive definite, and are $C^{1,1}([-a,a]\times \mathbb{R}; \mathbb{C}^{2\times 2})$. 
%In particular, for each $\vec{x}\in \mathbb{R}^2_a$, $\vec{p}\in \mathbb{C}^2$, 
%\begin{align*}
%	\begin{split}
%		\widetilde{\mathbb{T}}(\vec{x})\vec{p}\cdot 	\overline{\vec{p}}&\geq \chi_{\ell,2\delta}(y)\mathbb{A}_{11}^{-1}c_{\mathbb{T}}(\vec{x})|\vec{p}|^2+(1-\chi_{\ell,2\delta}(y))|\vec{p}|^2\\
%		&\geq 
%		\min(c_{\mathbb{T}}(\vec{x}), 1)(\chi_{\ell,2\delta}(y)|\vec{p}|^2+(1-\chi_{\ell,2\delta}(y))|\vec{p}|^2)=
%		\min(\min_{\vec{x}}c_{\mathbb{T}}(\vec{x}),1)|\vec{p}|^2,
%	\end{split}
%\end{align*}
%where in the last inequality we used $\chi_{\ell,2\delta}\in [0, 1]$.
Moreover, $\widetilde{\mathbb{D}}_{11}(\vec{x})=1$ for all $\vec{x}\in \mathbb{R}^2_a$. As we will see further, this will allow us to avoid appearance of  $\|F^{\nu}_{\delta}\|_{\mathbb{R}^2_a}\|\partial_x \mathcal{J}U^{\nu}_{\delta}\|_{\mathbb{R}^2_a}$ in the right-hand side, which would have prevented us from obtaining the sharp bound \eqref{eq:lemboundnu}. 
%
%Next, for $(x,y): \, |y|>\ell+2\delta$, the matrix $\widetilde{\mathbb{P}}$ becomes constant (actually, identity). Finally, it is straightforward to see that $\widetilde{\mathbb{P}}$ preserve its hermitian and positive definiteness properties: e.g. for $\mathbb{P}=\mathbb{T}$, for each $y\in \mathbb{R}$, $\vec{p}\in \mathbb{C}^2$, 

Let us now rewrite the problem satisfied by $U^{\nu}_{\delta}$ \eqref{eq:Unudelta}. We start by remarking that, as  $\operatorname{supp}U^{\nu}_{\delta}, \operatorname{supp}F^{\nu}_{\delta}\subseteq\operatorname{supp}\chi_{\ell,\delta}$, and  $\chi_{\ell,2\delta}=1$ on $\operatorname{supp}\chi_{\ell,\delta}$, the matrices $\mathbb{A}$ and $\mathbb{T}$ in \eqref{eq:Unudelta} can be replaced by $\mathbb{A}_{11}\widetilde{\mathbb{A}}, \, \mathbb{A}_{11}\widetilde{\mathbb{T}}$:
\begin{align*}
	\operatorname{div}\left(\left(x\mathbb{A}+i\nu\mathbb{T}\right)\nabla U^{\nu}_{\delta}\right)&=	\operatorname{div}\left(\mathbb{A}_{11}\left(x\widetilde{\mathbb{A}}+i\nu\widetilde{\mathbb{T}}\right)\nabla U^{\nu}_{\delta}\right)=\nabla \mathbb{A}_{11}\cdot \left(x\widetilde{\mathbb{A}}+i\nu\widetilde{\mathbb{T}}\right)\nabla U^{\nu}_{\delta}+	\mathbb{A}_{11}\operatorname{div}\left(\left(x\widetilde{\mathbb{A}}+i\nu\widetilde{\mathbb{T}}\right)\nabla U^{\nu}_{\delta}\right).
\end{align*}
The above can be rewritten as 
\begin{align}
	\label{eq:neweq}
	\operatorname{div}\left(\left(x\widetilde{\mathbb{A}}+i\nu\widetilde{\mathbb{T}}\right)\nabla U^{\nu}_{\delta}\right)=\widetilde{F}^{\nu}_{\delta},\quad \text{ in }\mathbb{R}^2_a,
\end{align}
where the right-hand side $\widetilde{F}^{\nu}_{\delta}$ satisfies the following bound, with some $C_j>0$, $j=1,2$, independent of $\nu>0$:
\begin{align}
	\label{eq:fnudelta_bound}
	\|\widetilde{F}^{\nu}_{\delta}\|_{L^2(\mathbb{R}^2_a)}\leq C_1(\|x\nabla U^{\nu}_{\delta}\|_{L^2(\mathbb{R}^2_a)}+\nu\|\nabla U^{\nu}_{\delta}\|_{L^2(\mathbb{R}^2_a)})+\|F^{\nu}_{\delta}\|_{L^2(\mathbb{R}^2_a)}\leq C_2\|f\|_{L^2(\Omega)}, 
\end{align}
with the latter bounds following from \eqref{eq:stabU} and \eqref{eq:stabFdelta}. Moreover, $\operatorname{supp}\widetilde{F}^{\nu}_{\delta}\subseteq {\operatorname{supp}}\chi_{\ell,2\delta}.$	


We test \eqref{eq:neweq} with $\mathcal{J}^2U^{\nu}_{\delta}$, which belongs to $H^1(\mathbb{R}^2_a)$, because of the identity \eqref{eq:Ju}  and the elliptic regularity for $U^{\nu}_{\delta}$ as stated in Lemma \ref{lem:ellreg}; moreover, $\mathcal{J}^2U^{\nu}_{\delta}\in H^1_0(\mathbb{R}^2)$ according to \eqref{eq:zero_trace}. Integrating by parts, we get
\begin{align*}
	\int_{\mathbb{R}^2_a}(x\widetilde{\mathbb{A}}&+i\nu\widetilde{\mathbb{T}})\nabla U^{\nu}_{\delta}\cdot \overline{\mathcal{J}^2\nabla U^{\nu}_{\delta}}=-\int_{\mathbb{R}^2_a}\widetilde{F}^{\nu}_{\delta}\overline{\mathcal{J}^2 U^{\nu}_{\delta}}.
\end{align*}
With Lemma \ref{lem:plancherel} on the symmetry of $\mathcal{J}$, we obtain the new identity:
\begin{align}
	\label{eq:ident}
	\begin{split}
		\int_{\mathbb{R}^2_a}(x\widetilde{\mathbb{A}}&+i\nu\widetilde{\mathbb{T}}) \mathcal{J}\nabla U^{\nu}_{\delta}\cdot \overline{\mathcal{J}\nabla U^{\nu}_{\delta}}+
		\overbrace{\int_{\mathbb{R}^2_a}[\mathcal{J}\mathbb{I},(x\widetilde{\mathbb{A}}+i\nu\widetilde{\mathbb{T}})]\nabla U^{\nu}_{\delta}\cdot \overline{\nabla \mathcal{J}U^{\nu}_{\delta}}}^{\mathcal{S}}=-\int_{\mathbb{R}^2_a}\widetilde{F}^{\nu}_{\delta}\overline{\mathcal{J}^2 U^{\nu}_{\delta}}.	
	\end{split}	
\end{align}
Taking the imaginary part of \eqref{eq:ident}, and using the fact that $\widetilde{\mathbb{A}}$ is a Hermitian matrix, and $\widetilde{\mathbb{T}}$ is Hermitian positive definite, yields the following inequality, with some $C>0$,
\begin{align}
	\label{eq:fbound}
	\nu\|\mathcal{J}\nabla U^{\nu}_{\delta}\|_{L^2(\mathbb{R}^2_a)}^2\leq C\left( |\Im \mathcal{S}|_{L^2(\mathbb{R}^2_a)}+\|\widetilde{F}^{\nu}_{\delta}\|_{L^2(\mathbb{R}^2_a)}\|\mathcal{J}^2 U^{\nu}_{\delta}\|_{L^2(\mathbb{R}^2_a)}\right).
\end{align}
The first term $|\Im\mathcal{S}|$ in the right-hand side of the above needs to be treated with care, since a na\"ive bound using Lemma \ref{lem:commutator} and \eqref{eq:stabU} would allow to replace  in \eqref{eq:fbound} $|\Im S|$ by  $\|F^{\nu}_{\delta}\|_{L^2(\mathbb{R}^2_a)}\|\nabla\mathcal{J}U^{\nu}_{\delta}\|_{L^2(\mathbb{R}^2_a)}$, which would not yield \eqref{eq:lemboundnu} because of the loss of a half a power of $\nu$.  

\textbf{Bounding $\Im\mathcal{S}$. }		
We rewrite $
	\mathcal{S}=\sum_{j}\mathcal{S}_j, $
with $\mathcal{S}_1, \mathcal{S}_2, \mathcal{S}_3$ defined as follows. With $\vec{v}=\nabla U^{\nu}_{\delta}$, and recalling that  $\widetilde{\mathbb{A}}=\widetilde{\mathbb{D}}+\widetilde{\mathbb{H}}$, we have that  
\begin{align}
	\label{eq:s123}
	\mathcal{S}_1=\int_{\mathbb{R}^2_a}x[\mathcal{J}\mathbb{I}, \widetilde{\mathbb{D}}]\vec{v}\cdot \overline{\mathcal{J}\vec{v}},\qquad 
	\mathcal{S}_2=\int_{\mathbb{R}^2_a}x[\mathcal{J}\mathbb{I}, \widetilde{\mathbb{H}}]\vec{v}\cdot \overline{\mathcal{J}\vec{v}},\qquad 			\mathcal{S}_3=i\nu \int_{\mathbb{R}^2_a}x[\mathcal{J}\mathbb{I}, \widetilde{\mathbb{T}}]\vec{v}\cdot \overline{\mathcal{J}\vec{v}}.
\end{align} 
\textbf{A bound on $\Im \mathcal{S}_1$. }We have, since $\widetilde{\mathbb{D}}_{11}=1$, see  \eqref{eq:ptilde}, and $\widetilde{\mathbb{D}}$ is a diagonal matrix:
\begin{align}
	\label{eq:I1}
	\Im\mathcal{S}_1=\Im \int_{\mathbb{R}^2_a}x[\mathcal{J}\mathbb{I}, \widetilde{\mathbb{D}}]\vec{v}\cdot \overline{\mathcal{J}\vec{v}}&=	\Im \int_{\mathbb{R}^2_a}x[\mathcal{J},\widetilde{\mathbb{D}}_{22}]v_2\cdot \overline{\mathcal{J}v_2}=\Im \int_{\mathbb{R}^2_a}x[\mathcal{J},\mathbb{E}_{\mathbb{D},22}]v_2\cdot \overline{\mathcal{J}v_2}.
\end{align}
Next, we make use of Lemma \ref{lem:self_adjoint_commutator}, with $\mathcal{A}\vec{v}=x\mathbb{E}_{\mathbb{D},22}v_2$ (we recall that $\mathbb{A}$ is Hermitian, thus $\mathbb{D}$, $\widetilde{\mathbb{D}}$, ${\mathbb{E}}_{\mathbb{D},22}$ are real-valued). This yields
\begin{align*}
	2i\Im \mathcal{S}_1&=\int_{\mathbb{R}^2_a}x[\mathcal{J}^2, \mathbb{E}_{\mathbb{D},22}]v_2\overline{v_2}.
\end{align*}
Next, we employ Lemma \ref{lem:commutator} with compactly supported  $\beta_0=\mathbb{E}_{\mathbb{D},22}=(\mathbb{A}_{22}/\mathbb{A}_{11}-1)\chi_{\ell,2\delta}$, cf. Remark \ref{rem:regularity}, and use the regularity Assumption \ref{assump:matrices}, which results in 
\begin{align*}
	|\Im \mathcal{S}_1|\leq C \|xv_2\|_{L^2(\mathbb{R}^2_a)}\|v_2\|_{L^2(\mathbb{R}^2_a)}.
\end{align*}
Recalling that $v_2=\partial_y U^{\nu}_{\delta}$, the bound \eqref{eq:stabU}, namely $\|x\nabla U^{\nu}_{\delta}\|_{L^2(\mathbb{R}^2_a)}\lesssim \|f\|_{L^2(\Omega)}$, and the bound \eqref{eq:important_mappings} on $\|\partial_y U^{\nu}_{\delta}\|_{L^2(\mathbb{R}^2_a)}$, we conclude that 
\begin{align}
	\label{eq:boundI1}
	\left|\Im \mathcal{S}_1\right|\leq C (\|u^{\nu}\|_{L^2(\Omega)}+\|\partial_y u^{\nu}\|_{L^2(\Omega)})\|f\|_{L^2(\Omega)}.
\end{align}
Remark that it is due to our rewriting \eqref{eq:neweq} that $\widetilde{\mathbb{D}}=1$, and we do not have terms of the type $\|x\partial_x U^{\nu}_{\delta}\|\|\partial_x U^{\nu}_{\delta}\|$ occurring in the bound on $\Im \mathcal{S}_1$.\\
\textbf{A bound on $\Im \mathcal{S}_2$. }
Let us consider
\begin{align}
	\label{eq:I2}
	\Im  \mathcal{S}_2=\operatorname{Im}\int_{\mathbb{R}^2_a}x[\mathcal{J}\mathbb{I}, \widetilde{\mathbb{H}}]\vec{v}\cdot\overline{\mathcal{J}\vec{v}}=\frac{1}{2i}
	\left(x[\mathcal{J}^2\mathbb{I}, \widetilde{\mathbb{H}}]\vec{v}, \vec{v}\right),
\end{align}
where the last identity follows again by Lemma \ref{lem:self_adjoint_commutator} with $\mathcal{A}\vec{v}=\widetilde{\mathbb{H}}\vec{v}$,  and using the fact that $\widetilde{\mathbb{H}}$ is Hermitian. 
Since $\operatorname{diag}\widetilde{\mathbb{H}}=0$, and using the decomposition \eqref{eq:ptilde}, we rewrite the above as follows:
\begin{align*}
	|\Im \mathcal{S}_2|&\lesssim \|x[\mathcal{J}^2, \widetilde{\mathbb{H}}_{21}]v_1\|_{L^2(\mathbb{R}^2_a)}\|v_2\|_{L^2(\mathbb{R}^2_a)}+\|[\mathcal{J}^2, \widetilde{\mathbb{H}}_{12}]v_2\|_{L^2(\mathbb{R}^2_a)}\|xv_1\|_{L^2(\mathbb{R}^2_a)}\\
	&=\|x[\mathcal{J}^2, \mathbb{E}_{\mathbb{H},21}]v_1\|_{L^2(\mathbb{R}^2_a)}\|v_2\|_{L^2(\mathbb{R}^2_a)}+\|[\mathcal{J}^2, \mathbb{E}_{\mathbb{H},12}]v_2\|_{L^2(\mathbb{R}^2_a)}\|xv_1\|_{L^2(\mathbb{R}^2_a)}.
\end{align*}
By the repeated application of Lemma \ref{lem:commutator}, first with $\beta_0=\mathbb{E}_{\mathbb{H},21}$, next with $\beta_0=\mathbb{E}_{\mathbb{H},12}$, and using $\vec{v}=\nabla U^{\nu}_{\delta}$,  
\begin{align}
	\label{eq:boundI2}
	|\Im \mathcal{S}_2|\lesssim \|x\partial_x U^{\nu}_{\delta}\|_{L^2(\mathbb{R}^2_a)}\|\partial_y U^{\nu}_{\delta}\|_{L^2(\mathbb{R}^2_a)}\lesssim \|f\|_{L^2(\Omega)}(\|u^{\nu}\|_{L^2(\Omega)}+\|\partial_y u^{\nu}\|_{L^2(\Omega)}),
\end{align}	
with the last bound obtained like in \eqref{eq:boundI1}.\\
\textbf{A bound on $\Im \mathcal{S}_3$. }It remains to consider the remaining term, namely,  
\begin{align}
	\label{eq:I3}
	\Im\mathcal{S}_3=\nu\Re \int_{\mathbb{R}_a^2}[\mathcal{J}\mathbb{I}, \mathbb{\widetilde{T}}]\vec{v}\cdot\overline{\mathcal{J}\vec{v}}=\nu\Re \int_{\mathbb{R}_a^2}[\mathcal{J}\mathbb{I}, \mathbb{E}_{\mathbb{T}}]\vec{v}\cdot\overline{\mathcal{J}\vec{v}}.
\end{align}
With the Cauchy-Schwarz inequality, we obtain that 
\begin{align*}
	|\Im \mathcal{S}_3|\leq \nu \|[\mathcal{J}\mathbb{I}, \mathbb{E}_{\mathbb{T}}]\vec{v}\|_{L^2(\mathbb{R}^2_a)}\|{\mathcal{J}\vec{v}}\|_{L^2(\mathbb{R}^2_a)}.
\end{align*}
Next we use Lemma \ref{lem:commutator} (where $\beta_0=\mathbb{E}_{\mathbb{T},ij},\; i,j\in\{1,2\}$), which gives, together with $\vec{v}=\nabla U^{\nu}_{\delta}$,  
\begin{align}
	\label{eq:boundI3}
	|\Im \mathcal{S}_3|\lesssim \nu \|\nabla U^{\nu}_{\delta}\|_{L^2(\mathbb{R}^2_a)}\|\mathcal{J}\nabla U^{\nu}_{\delta}\|_{L^2(\mathbb{R}^2_a)}\overset{\eqref{eq:stabU}}{\lesssim} \nu^{1/2}\|f\|_{L^2(\Omega)}\|\mathcal{J}\nabla U^{\nu}_{\delta}\|_{L^2(\mathbb{R}^2_a)}.
\end{align}	
\textbf{The final bound on $\Im \mathcal{S}$. }	Gathering \eqref{eq:boundI1}, \eqref{eq:boundI2}, \eqref{eq:boundI3} into $\mathcal{S}=\sum\limits_{j=1}^3\mathcal{S}_j$, we conclude that 
\begin{align}
	\label{eq:boundImI}
	|\Im \mathcal{S}|\leq \nu^{1/2}\|f\|_{L^2(\Omega)}\|\mathcal{J}\nabla U^{\nu}_{\delta}\|_{L^2(\mathbb{R}^2_a)}+\|f\|_{L^2(\Omega)}(\|u^{\nu}\|_{L^2(\Omega)}+\|\partial_y u^{\nu}\|_{L^2(\Omega)}).
\end{align}
\textbf{Bounding $\mathcal{J}\nabla U^{\nu}_{\delta}$. }We get back to \eqref{eq:fbound}. We employ  \eqref{eq:boundImI}, the bound for $\widetilde{F}^{\nu}_{\delta}$ in \eqref{eq:fnudelta_bound}, as well as \eqref{eq:Ju}, namely, $\|\mathcal{J}^2U^{\nu}_{\delta}\|_{L^2(\mathbb{R}^2_a)}\leq \left(\|U^{\nu}_{\delta}\|_{L^2(\mathbb{R}^2_a)}+\|\partial_y U^{\nu}_{\delta}\|_{L^2(\mathbb{R}^2_a)}\right)$. This results in the following inequality:
\begin{align*}
	\nu\|\mathcal{J}\nabla U^{\nu}_{\delta}\|^2_{L^2(\mathbb{R}^2_a)}&\lesssim \nu^{1/2}\|f\|_{L^2(\Omega)}\|\mathcal{J}\nabla U^{\nu}_{\delta}\|_{L^2(\mathbb{R}^2_a)}+\|f\|_{L^2(\Omega)}\left(\|\partial_y u^{\nu}\|_{L^2(\Omega)}+\|u^{\nu}\|_{L^2(\Omega)}\right)\\
	&+\|f\|_{L^2(\Omega)}(\|U^{\nu}_{\delta}\|_{L^2(\mathbb{R}^2_a)}+\|\partial_y U^{\nu}_{\delta}\|_{L^2(\mathbb{R}^2_a)}).
\end{align*}
Next, we use \eqref{eq:important_mappings}, which allows to replace $U^{\nu}_{\delta}$ by $u^{\nu}$ in the last term:
\begin{align*}
	\nu\|\mathcal{J}\nabla U^{\nu}_{\delta}\|^2_{L^2(\mathbb{R}^2_a)}&\lesssim \nu^{1/2}\|f\|_{L^2(\Omega)}\|\mathcal{J}\nabla U^{\nu}_{\delta}\|_{L^2(\mathbb{R}^2_a)}+\|f\|_{L^2(\Omega)}\left(\|\partial_y u^{\nu}\|_{L^2(\Omega)}+\|u^{\nu}\|_{L^2(\Omega)}\right).
\end{align*}
Finally, with \eqref{eq:est_main_v2}, we deduce that the following inequality holds true with some $C>0$ independent of $\nu$, uniformly in $0<\nu<1$:
\begin{align*}
	\nu\|\mathcal{J}\nabla U^{\nu}_{\delta}\|^2_{L^2(\mathbb{R}^2_a)}\leq C( \|f\|_{L^2(\Omega)}\times \nu^{1/2}\|\mathcal{J}\nabla U^{\nu}_{\delta}\|_{L^2(\mathbb{R}^2_a)}+\|f\|_{L^2(\Omega)}\|\partial_y u^{\nu}\|_{L^2(\Omega)}+\|f\|^2_{L^2(\Omega)})
\end{align*}
and the desired bound follows by applying the Young's inequality to the above. 
\end{proof}

\end{document}
